% figuras/tikz/cap02-rigida-vs-difusa.tex
% Contraste conceptual: segmentación rígida (pertenencia binaria, k-means)
% vs. segmentación difusa (pertenencia parcial, fuzzy c-means).
\begin{tikzpicture}[
    seg/.style={circle, draw=guindaIPN, thick, fill=guindaIPNClara, minimum size=2.6cm, font=\scriptsize, align=center},
    segb/.style={circle, draw=bronceIPN, thick, fill=bronceIPNClaro, minimum size=2.6cm, font=\scriptsize, align=center},
    punto/.style={circle, fill=black, minimum size=3pt, inner sep=0pt},
    titulo/.style={font=\small\bfseries}
  ]
  % Panel izquierdo: segmentación rígida
  \node[titulo] at (0,2.2) {Rígida ($k$-means)};
  \node[seg] (A1) at (-1.1,0) {Segmento A};
  \node[segb] (B1) at (1.5,0) {Segmento B};
  \node[punto] (p1) at (0.15,0.15) {};
  \node[font=\tiny, align=center] at (0.15,-0.55) {consumidor $i$\\ pertenece solo a A};

  % Panel derecho: segmentación difusa
  \begin{scope}[xshift=6.5cm]
    \node[titulo] at (0,2.2) {Difusa (\textit{fuzzy c-means})};
    \node[seg, opacity=0.75] (A2) at (-0.9,0) {Segmento A};
    \node[segb, opacity=0.75] (B2) at (0.9,0) {Segmento B};
    \node[punto] (p2) at (0,0.1) {};
    \node[font=\tiny, align=center] at (0,-0.85) {consumidor $i$: $u_{iA}=0.63$,\\ $u_{iB}=0.37$};
  \end{scope}
\end{tikzpicture}
